% Partie : logique intuitionniste
% Chapitre : propositions comme types
% Section : termes de preuve

\documentclass[../../../include/open-logic-section]{subfiles}

\begin{document}

\olfileid[fr]{int}{pty}{ter}

\olsection{Termes de preuve}

Dans \Log{N1i}, nous étiquetons les hypothèses par des variables,
par exemple $!A^x$, pour indiquer leur décharge.
Dans \Log{N2i}, le contexte~$\Gamma$ à gauche de chaque
séquent $\Gamma \Sequent !A$ est constitué de couples
variable--formule. Chaque séquent d'une preuve indique donc
la formule déjà prouvée, $!A$, et les hypothèses ouvertes
à ce stade, $\Gamma$.
Nous nous limitons au fragment engendré par les atomes,
l'absurde, l'implication, la conjonction et la disjonction.
Les termes ci-dessous correspondent aux règles de ce fragment.
Ajoutons à chaque séquent une information sur la manière
dont $!A$ a été prouvée. Nous obtenons un système
\emph{étiqueté par des termes}, \Log{tN2i}, dont les séquents
ont la forme $\Gamma \Sequent N : !A$.
Comme auparavant, $\Gamma$ est un ensemble de couples
$x: !B$, et $!A$ est la formule que la sous-preuve
achevée par ce séquent établit à partir des hypothèses
de~$\Gamma$. Le terme $N$ encode cette sous-preuve.

Les termes attribués aux séquents sont construits à partir
des variables $x$, $y$, \dots, à l'aide de symboles de fonction
qui encodent les inférences effectuées.
Chaque règle d'inférence reçoit son propre symbole.
Ceux associés aux règles d'introduction sont des
« constructeurs » ; ceux associés aux règles d'élimination,
des « destructeurs ». Le nombre d'arguments d'un symbole
correspond au nombre de prémisses de la règle.
Ainsi, pour \Intro{\land} et \Elim{\land}[1], on pourrait
employer respectivement $c^\land$, d'arité deux,
et $d^\land_1$, d'arité un.
Les règles correspondantes de \Log{tN2i} seraient :
\[
\Axiom$\Gamma_1 \fCenter N: !A$
\Axiom$\Gamma_2 \fCenter M: !B$
\RightLabel{\Intro{\land}}
\BinaryInf$\Gamma_1, \Gamma_2 \fCenter c^\land(N, M) : !A \land !B $
\DisplayProof
\quad
\Axiom$\Gamma \fCenter N: !A \land !B$
\RightLabel{\Elim{\land}[1]}
\UnaryInf$\Gamma \fCenter d^\land_1(N) : !A$
\DisplayProof
\]
Lorsqu'une règle décharge une hypothèse, le terme de preuve
doit également l'indiquer : une formule étiquetée~$x:!A$
figure dans une prémisse mais n'est plus une hypothèse
ouverte dans la conclusion. Les symboles associés à ces
règles \emph{lient} les variables correspondantes.
Par exemple, \Intro{\lif} passe de
$x:!A, \Gamma \Sequent !B$ à $\Gamma \Sequent !A \lif !B$.
Le constructeur $c^{\lif}$ prend un argument et indique
que la variable $x$ devient liée : c'est ainsi que le terme
encode la décharge de l'hypothèse étiquetée par~$x$.
Si le terme de preuve de la prémisse est~$N$, celui de
la conclusion peut s'écrire $c^{\lif}([x]N)$ :
\[
\Axiom$x:!A, \Gamma \fCenter N: !B$
\RightLabel{\Intro{\lif}}
\UnaryInf$\Gamma \fCenter c^{\lif}([x]N) : !A \lif !B$
\DisplayProof
\]

Pour reconstruire toute la preuve, une fois connu le
contexte des variables libres, certaines annotations
supplémentaires sont nécessaires. Lorsque les formules
des prémisses ne déterminent pas entièrement celle de
la conclusion, il faut enregistrer l'information manquante
dans le terme. C'est le cas de \Intro{\lif}, pour laquelle
nous écrivons $c^{\lif}([x:!A]N)$.
C'est aussi le cas de \Intro{\lor} et de~\FalseInt :
\[
\Axiom$\Gamma \fCenter N:!A$
\RightLabel{\Intro{\lor}[1]}
\UnaryInf$\Gamma \fCenter c^{\lor{!B}}_1(N):!A \lor !B$
\DisplayProof
\quad
\Axiom$\Gamma \fCenter N:\lfalse$
\RightLabel{\FalseInt}
\UnaryInf$\Gamma \fCenter d^\lfalse_{!A}(N):!A$
\DisplayProof
\]

On peut ainsi définir une présentation séquentielle de
la déduction naturelle dont tous les séquents ont la forme
$\Gamma \Sequent N : !A$, où $N$ est un \emph{terme de preuve}.

Ces termes sont étroitement liés aux termes du
\emph{lambda-calcul typé}. Pour faire apparaître ce lien,
nous utiliserons les notations usuelles du lambda-calcul
plutôt que les symboles génériques précédents.
Nous écrirons $\lambd[\typeof{x}{!A}][N]$ à la place de
$c^{\lif}([x:!A]N)$, $(NM)$ à la place de $d^{\lif}(N,M)$,
$\pair{N}{M}$ à la place de $c^\land(N,M)$, et $\proj{i}{N}$
à la place de~$d^\land_i(N)$.

\begin{defn}
Les termes de preuve sont engendrés inductivement par
les règles suivantes :
\begin{enumerate}
\item Toute variable $x$ est un terme de preuve (axiomes).
\item Si $x$ est une variable, $!A$ une formule et $N$
  un terme de preuve, alors $\lambd[\typeof{x}{!A}][N]$
  est un terme de preuve~(\Intro\lif).
\item Si $N$ et $M$ sont des termes de preuve, alors $(NM)$
  est un terme de preuve~(\Elim\lif).
\item Si $N$ et $M$ sont des termes de preuve, alors
  $\pair{N}{M}$ est un terme de preuve~(\Intro\land).
\item Si $N$ est un terme de preuve, alors $\proj{i}{N}$
  est un terme de preuve, pour $i$ égal à $1$ ou à~$2$
  ~(\Elim\land[i]).
\item Si $N$ est un terme de preuve et $!A$ une formule,
  alors $\inj[!A]{i}{N}$ est un terme de preuve,
  pour $i$ égal à $1$ ou à~$2$~(\Intro\lor[i]).
\item Si $M$, $N_1$ et $N_2$ sont des termes de preuve,
  et $x_1$, $x_2$ des variables, alors
  $\dcase{M}{x_1}{N_1}{x_2}{N_2}$ est un terme de preuve
  ~(\Elim\lor).
\item Si $N$ est un terme de preuve et $!A$ une formule,
  alors $\abort{!A}{N}$ est un terme de preuve~(\FalseInt).
\end{enumerate}
\end{defn}

Dans l'abstraction, la variable indiquée est liée dans le corps.
Dans l'analyse par cas, chaque variable de branche est liée
uniquement dans le terme de sa propre branche ; elle n'est
pas liée dans le terme analysé ni dans l'autre branche.
Les variables liées peuvent être renommées pour éviter
la capture et les conflits avec les variables du contexte.

Tout terme de preuve ainsi construit n'encode pas
nécessairement une preuve correcte.
Par exemple, $\pair{N}{M}$ encode une preuve dont la
dernière inférence est \Intro{\land}, et dont la formule
finale est donc une conjonction $!A \land !B$.
Elle ne peut pas servir de prémisse majeure à \Elim{\lif}.
Le terme $(\pair{N}{M}P)$ ne peut donc pas être le terme
d'une preuve correcte. Les termes \emph{corrects}
sont ceux qui respectent les règles de \Log{tN2ip},
données dans le \olref{tab:tN2ip}.

\subfile{rules-tN2}

\end{document}
